\Bezier curves are one of the most common representations for free-form
curves in computer graphics. The curves are named for Pierre \Bezier,
one of the people who were instrumental in their development. \Bezier
curves have an interesting history where they were concurrently
developed by several independent groups.

A \Bezier curve is a polynomial curve that approximates its control
points. The curves can be any degree of polynomial. A curve of degree
$d$ is controlled by $d+1$ control points. The curve interpolates its
first and last control points, and the shape is directly influenced by
the other points.

Often, complex shapes are made by connecting a number of \Bezier curves
of low degree. To avoid confusion, we refer to each polynomial piece
as a segment. In computer graphics, cubic ($d=3$) \Bezier curves are
the most commonly used for each segment. Many popular illustration
programs, such as Adobe Illustrator, and font representation schemes,
such as that used in Postscript, use cubic \Bezier segments.
\Bezier curves are extremely popular in computer graphics because they
are easy to control, have a number of useful properties, and there are
very efficient algorithms for working with them.

\Bezier curves are constructed such that:
\begin{itemize}
\item the curve interpolates the first and last control points,
      with $u=0$ and $1$ respectively;
\item the first derivative of the curve at its begining (end) is
      determined by the vector between the first and second (next
      to last and last) control points. The derivatives are given
      by the vectors between these points scaled by the degree of
      the curve;
\item higher derivatives at the begining (end) of the curve depend
         on the points at the begining (end) of the curve. The
         $n^{th}$ derivative depends on the first (last) $n+1$ points.
\end{itemize}

\begin{figure}
\centerline{\includegraphics{Figs/new-bezier-cubic.pdf}}
\caption{\label{fig:cubic-bezier}
	A cubic \Bezier segment is controlled by four points.
	It interpolates the first and last, and the begining and final
	derivatives are 1/3 of the vectors between the first two (or
	last two) points.
	}
\end{figure}


For example, consider the \Bezier curve of degree 3 as in
Figure~\ref{fig:cubic-bezier}. The curve has four ($d+1$) control
points. It begins at the first control point ($\mathbf{p_0}$) and ends
at the last ($\mathbf{p_1}$). The first derivative at the begining is
proportional to the vector between the first and second control points
$(\mathbf{p_1}-\mathbf{p_0})$. Specifically,
$\mathbf{f'}(0)=3(\mathbf{p_1}-\mathbf{p_0})$. Similarly, the first
derivative at the end of the curve is given by
$\mathbf{f'}(1)=3(\mathbf{p_3}-\mathbf{p_2})$. The second derivative
at the begining of the curve can be determined from control points
$\mathbf{p_0}$, $\mathbf{p_1}$ and $\mathbf{p_2}$.

Using the facts about Bezier cubics in the preceeding paragraph, we
can use the methods of Section \ref{sec:cubics} to create a parametric
function for them. The definitions of the begining and end
interpolation and derivatives give:
\[
\begin{array}{rcl}
	 \mathbf{p_0} = & \mathbf{f}(0) &
		=   \mathbf{a_3} 0^3 + \mathbf{a_2} 0^2 + \mathbf{a_1} 0 + \mathbf{a_0}\\
	 \mathbf{p_3} = & \mathbf{f}(1)&
		=   \mathbf{a_3} 1^3 + \mathbf{a_2} 1^2 + \mathbf{a_1} u + \mathbf{a_0} \\
	 3(\mathbf{p_1}-\mathbf{p_0}) = & \mathbf{f'}(0) &
		= 3 \mathbf{a_3} 0^2 + 2 \mathbf{a_2} 0 + \mathbf{a_1}     \\
	 3(\mathbf{p_3}-\mathbf{p_2}) =&  \mathbf{f'}(1) &
		= 3 \mathbf{a_3} 1^2 + 2 \mathbf{a_2} 1 + \mathbf{a_1}.
\end{array}
\]
This can be solved for the basis matrix
\[
	\mathbf{B} = \mathbf{C}^{-1} =
	\left[\begin{array}{rrrr}
	 1&  0&  0& 0\\
	-3&  3&  0& 0\\
	 3& -6&  3& 0\\
	-1&  3& -3& 1\\
	\end{array}\right],
\]
   and then written as
\[
       \mathbf{f}(u) = (1-3u+3u^2-u^3) \mathbf{p_0}
		     + (3u-6u^2+3u^3) \mathbf{p_1}
		     + (3u^2-3u^3) \mathbf{p_2}
		     + (u^3) \mathbf{p_3},
\]
   or
\[
       \mathbf{f}(u) = \displaystyle\sum_{i=0}^{d} b_{i,3}\mathbf{p_i}
\]
where the $b_{i,3}$ are the \Bezier blending functions of degree 3:
\[
	\begin{array}{rl}
       b_{0,3} = & (1-u)^3 \\
       b_{1,3} = & 3u(1-u)^2\\
       b_{2,3} = & 3u^2(1-u)\\
       b_{3,3} = & u^3.
	\end{array}
\]

Fortunately, the blending functions for Bezier curves have a special
form that works for all degrees. These functions are known as the
{\em Bernstein Basis Polynomials} and have the general form
\[
b_{k,n}(u)\ =\ C(n,k)\ u^k\ (1-u)^{(n-k)},
\]
